\section{Por qué esta es la «última entrevista»}

Lo particular de este episodio es que, cuando se grabó, la adquisición todavía no ocurría, y cuando se publicó, el desenlace ya se conocía. Por eso el lector puede ver un Manus que la narrativa de la adquisición aún no reescribía: todavía habla de su propio crecimiento, de sus riesgos, de la complejidad, del mercado de Agent y de la próxima década, en lugar de buscar después una explicación para la operación.

\subsection{Lectura posterior y lectura del momento}

Leído en su momento, el foco de Peak está en el producto y en la empresa: cómo crecer, cómo evitar la complejidad, cómo enfrentar a las grandes empresas tecnológicas y cómo llevar los Agent a más personas. Leído después, la adquisición por parte de Meta le añade a este contenido otra capa de significado: una empresa de Agent que todavía insiste en la simplicidad y en su identidad propia, ¿podrá conservar esa identidad una vez integrada en una plataforma más grande? Esta es también la pregunta de este episodio que más vale la pena seguir observando.

\begin{knowledgebox}{El valor de la última entrevista}
El valor de la última entrevista no está en que haya anticipado la adquisición, sino en que registra los juicios de producto originales, previos a ella. Nos permite comparar: una vez que el equipo entra en una plataforma grande, qué juicios se conservan y cuáles reescribirán la organización y la estrategia.
\end{knowledgebox}

\subsection{La tensión latente entre Manus y Meta}

Meta tiene modelos, grafo social, infraestructura y experiencia en productos globales; Manus tiene una idea más clara de lo que es un producto de Agent en la mente del usuario y retroalimentación de sus primeros usuarios. Las posibilidades de combinar ambas cosas son grandes, pero la tensión también es evidente: una plataforma grande tiende a ampliar funciones, integrar su ecosistema y atender más puntos de entrada; lo que Manus más teme es precisamente perder su identidad y volverse complejo.

\begin{warningbox}{El riesgo central después de la adquisición}
Que una plataforma grande adquiera un producto no lo hace más fuerte de forma automática. El verdadero reto es si, con más recursos, Manus podrá seguir siendo contenido y seguir dejando claro al usuario qué problema resuelve.
\end{warningbox}

\subsection{Resumen de la sección}

Como última entrevista, el valor de este episodio está en que conserva la descripción que Manus hacía de sí mismo antes de ser absorbido por una plataforma grande. Cualquier evaluación posterior del éxito o el fracaso de la adquisición debería volver a estos juicios originales.
